\documentclass[12pt,a4paper]{article}

\usepackage[utf8]{inputenc}
\usepackage[T1]{fontenc}
\usepackage[spanish,es-nodecimaldot]{babel}
\usepackage{amsmath,amssymb,amsfonts,mathtools}
\usepackage{booktabs}
\usepackage{listings}
\usepackage{xcolor}
\usepackage{geometry}
\usepackage{enumitem}
\usepackage{hyperref}
\usepackage{graphicx}
\usepackage{float}
\usepackage{caption}
\usepackage{array}

\geometry{top=2.5cm, bottom=2.5cm, left=2.5cm, right=2.5cm}

% ── Colores ──────────────────────────────────────────────────
\definecolor{codeblue}{RGB}{0,70,180}
\definecolor{codegray}{RGB}{120,120,120}
\definecolor{codegreen}{RGB}{0,120,0}
\definecolor{codepurple}{RGB}{130,0,130}
\definecolor{backcolor}{RGB}{248,248,250}
\definecolor{bordercolor}{RGB}{200,200,210}

% ── Estilo listados CUDA/C++ ──────────────────────────────────
\lstdefinestyle{cuda}{
  language=C++,
  backgroundcolor=\color{backcolor},
  basicstyle=\ttfamily\footnotesize,
  keywordstyle=\color{codeblue}\bfseries,
  commentstyle=\color{codegreen}\itshape,
  stringstyle=\color{codepurple},
  numbers=left,
  numberstyle=\tiny\color{codegray},
  stepnumber=1,
  numbersep=8pt,
  frame=single,
  rulecolor=\color{bordercolor},
  breaklines=true,
  breakatwhitespace=false,
  tabsize=4,
  showstringspaces=false,
  morekeywords={__global__,__device__,__shared__,__host__,
                threadIdx,blockIdx,blockDim,gridDim,
                cudaMalloc,cudaFree,cudaMemcpy,cudaMemset,
                cudaMemcpyHostToDevice,cudaMemcpyDeviceToHost,
                cudaDeviceProp,cudaGetDeviceProperties,
                cublasHandle_t,cublasCreate,cublasDestroy,
                cublasSgemv,CUBLAS_OP_N,
                uint8_t,uint32_t,
                vector,string,unique_ptr,make_unique,
                syncthreads,fmaxf,sqrtf,expf,logf},
  escapeinside={(*@}{@*)},
  captionpos=b,
}
\lstset{style=cuda}

\hypersetup{
  colorlinks=true,
  linkcolor=codeblue,
  urlcolor=codeblue,
  citecolor=codegreen,
}

% ── Macros matemáticas ────────────────────────────────────────
\newcommand{\R}{\mathbb{R}}
\DeclareMathOperator*{\argmax}{argmax}
\newcommand{\todo}[1]{\textit{\textcolor{codegray}{[Pendiente: #1]}}}

% ─────────────────────────────────────────────────────────────
\title{%
  \textbf{Clasificación de Imágenes con CNN + MLP Acelerado en CUDA}\\[0.5em]
  \large Implementación desde cero de un pipeline\\
  Red Neuronal Convolucional (CPU) + Perceptrón Multicapa CUDA (GPU)
}
\author{}
\date{}

\begin{document}

\maketitle
\tableofcontents
\newpage

% ============================================================
\section{Introducción}
% ============================================================

La clasificación de imágenes es una de las tareas centrales del aprendizaje automático moderno.
Aplicar un Perceptrón Multicapa (MLP) directamente sobre los píxeles crudos presenta tres
limitaciones estructurales fundamentales:

\begin{itemize}[noitemsep]
  \item \textbf{Ignora la estructura espacial:} los píxeles adyacentes son igualmente
    importantes para el modelo que píxeles ubicados en extremos opuestos de la imagen.
  \item \textbf{No tiene invarianza traslacional:} un dígito ``3'' desplazado dos píxeles
    a la derecha produce una representación completamente distinta en el vector de entrada.
  \item \textbf{Escala cuadrática de parámetros:} una imagen de $32\times32$ ya genera
    1\,024 entradas; con una primera capa de 512 neuronas el modelo tiene más de 500\,000
    parámetros solo en la primera capa.
\end{itemize}

Las Redes Neuronales Convolucionales (CNN) resuelven estas limitaciones explotando tres
principios:

\begin{enumerate}[noitemsep]
  \item \textbf{Localidad:} cada neurona procesa únicamente una región local
    (\textit{campo receptivo}) en lugar de toda la imagen.
  \item \textbf{Compartición de parámetros:} el mismo filtro se aplica en todas las
    posiciones espaciales, reduciendo drásticamente el número de pesos entrenables.
  \item \textbf{Equivarianza traslacional:} si la entrada se desplaza, las activaciones
    internas se desplazan del mismo modo, facilitando la detección de características
    en cualquier posición de la imagen.
\end{enumerate}

Este informe documenta la implementación completa de un pipeline
\textbf{CNN (CPU) + MLP CUDA (GPU)} desarrollado en C++17/CUDA sin frameworks de alto nivel.
La retropropagación fluye de forma continua desde la capa de salida del MLP hasta los kernels
de la primera capa convolucional, actualizando todos los parámetros del modelo de forma
coordinada.

La división CPU/GPU responde a un compromiso de diseño deliberado:
\begin{itemize}[noitemsep]
  \item Las capas convolucionales corren en \textbf{CPU en precisión \texttt{double} (64 bits)}
    para garantizar estabilidad numérica en los gradientes de los kernels, que son pequeños
    y sensibles al redondeo.
  \item El MLP corre en \textbf{GPU en precisión \texttt{float} (32 bits)} y se beneficia
    de la tasa de cómputo masivamente paralela de CUDA y la biblioteca cuBLAS para
    operaciones matriciales.
\end{itemize}

% ============================================================
\section{Fundamentos Teóricos}
% ============================================================

\subsection{La Operación de Convolución 2D}

Dado un tensor de entrada $\mathbf{X} \in \R^{C \times H \times W}$ (canales, alto, ancho)
y un banco de $F$ filtros $\mathbf{K} \in \R^{F \times C \times K_h \times K_w}$, la
pre-activación del mapa de características del filtro $f$ en la posición $(h', w')$
con paso \textit{stride} $s$ es:

\begin{equation}
  \mathbf{Z}[f,\, h',\, w'] \;=\; b_f \;+\;
  \sum_{c=0}^{C-1} \sum_{i=0}^{K_h-1} \sum_{j=0}^{K_w-1}
  \mathbf{X}_{\mathrm{pad}}[c,\; h'\!\cdot s + i,\; w'\!\cdot s + j]
  \cdot \mathbf{K}[f,\, c,\, i,\, j]
  \label{eq:conv}
\end{equation}

donde $b_f$ es el sesgo del filtro $f$ y $\mathbf{X}_{\mathrm{pad}}$ es el input con
relleno de ceros (\textit{zero-padding}). En la práctica del deep learning la operación
\eqref{eq:conv} es técnicamente una \textbf{correlación cruzada} (el kernel no se voltea
espacialmente), convención estándar adoptada en la implementación.

\subsubsection{Modos de Relleno: VALID y SAME}

El modo de relleno determina la relación entre las dimensiones del input y del output:

\begin{itemize}
  \item \textbf{VALID} (sin relleno): el kernel solo se aplica donde cabe completamente
    dentro del input.
    \[H' = \left\lfloor \frac{H - K_h}{s} \right\rfloor + 1\]

  \item \textbf{SAME}: se añaden ceros simétricamente para que la salida tenga
    $H' = \lceil H / s \rceil$, conservando las dimensiones espaciales cuando $s=1$.
    El relleno total necesario en el eje vertical es:
    \begin{equation}
      p_H = \max\!\left(0,\;
        \left\lceil \frac{H}{s} \right\rceil \cdot s - H + K_h - s
      \right)
      \label{eq:same_pad}
    \end{equation}
    Se distribuye $\lfloor p_H/2 \rfloor$ arriba y $\lceil p_H/2 \rceil$ abajo
    (análogo para el eje horizontal).
\end{itemize}

La implementación utiliza SAME con $s=1$ en ambas capas convolucionales, por lo que los
mapas de características conservan exactamente las dimensiones espaciales del input a
cada capa convolucional.

\subsubsection{Activación ReLU y Caché para Backprop}

Tras la suma ponderada se aplica ReLU elemento a elemento:
\begin{equation}
  \mathbf{A}[f, h', w'] = \max\!\left(0,\; \mathbf{Z}[f, h', w']\right)
\end{equation}

La capa guarda en caché el input original, el input con padding y la pre-activación
$\mathbf{Z}$ antes de ReLU, ya que los tres son necesarios durante la retropropagación.

\subsubsection{Inicialización He de los Kernels}

Los pesos del kernel se inicializan con la estrategia de He (Kaiming):
\begin{equation}
  w \sim \mathcal{N}\!\left(0,\;\; \frac{2}{C \cdot K_h \cdot K_w}\right)
\end{equation}
El factor 2 en el numerador compensa que ReLU anula en promedio la mitad de las
activaciones, manteniendo la varianza de la señal estable a través de las capas.

% ──────────────────────────────────────────────────────────────
\subsection{Capa de Pooling}

El pooling reduce las dimensiones espaciales del mapa de características mediante una
operación no lineal sobre ventanas locales, proporcionando invarianza a pequeñas
traslaciones y reduciendo el coste computacional de las capas siguientes.

Dada una ventana de $P \times P$ con paso $s_P$, las dimensiones de salida son:
\begin{equation}
  H'_{\mathrm{pool}} = \left\lfloor \frac{H - P}{s_P} \right\rfloor + 1
\end{equation}

\subsubsection{MAX Pooling}

Selecciona el valor máximo dentro de cada ventana:
\begin{equation}
  \mathbf{M}[c, h', w'] = \max_{\substack{0 \le i < P \\ 0 \le j < P}}
  \mathbf{A}[c,\; h' \cdot s_P + i,\; w' \cdot s_P + j]
\end{equation}

La implementación almacena en una \textbf{máscara de posiciones ganadoras} el índice
lineal $\texttt{mask\_}[c][h'][w'] = i \cdot P + j$ del elemento que produjo el máximo.
Esta máscara es la única información necesaria para el backward.

\subsubsection{AVERAGE Pooling}

Calcula el promedio de la ventana, distribuyendo el gradiente uniformemente durante
la retropropagación:
\begin{equation}
  \mathbf{M}[c, h', w'] = \frac{1}{P^2}
  \sum_{i=0}^{P-1} \sum_{j=0}^{P-1} \mathbf{A}[c,\; h' \cdot s_P + i,\; w' \cdot s_P + j]
\end{equation}

% ──────────────────────────────────────────────────────────────
\subsection{Capa de Aplanado (Flatten)}

La capa de aplanado convierte el tensor tridimensional de salida del pooling
$\mathbf{M} \in \R^{C \times H' \times W'}$ en un vector unidimensional:
\begin{equation}
  \mathbf{x}_{\mathrm{flat}}[\,c \cdot H' \cdot W' + h' \cdot W' + w'\,]
  = \mathbf{M}[c,\, h',\, w']
\end{equation}

Es una operación puramente topológica, sin parámetros entrenables. El vector resultante
$\mathbf{x}_{\mathrm{flat}} \in \R^{C \cdot H' \cdot W'}$ constituye la entrada al MLP
y cierra el puente entre la \textbf{extracción de características} (parte CNN) y la
\textbf{clasificación} (parte MLP).

% ──────────────────────────────────────────────────────────────
\subsection{Retropropagación en la CNN}

La retropropagación aplica la regla de la cadena a través de cada capa en orden inverso.
El gradiente $\partial \mathcal{L}/\partial \mathbf{x}_{\mathrm{flat}}$ llega desde el
MLP en la GPU y debe propagarse hacia atrás por Flatten, Pooling y Convolución.

\subsubsection{Gradiente a través de Flatten}

El backward de flatten remodela el vector gradiente de vuelta a la forma del tensor:
\begin{equation}
  \frac{\partial \mathcal{L}}{\partial \mathbf{M}[c, h', w']}
  = \frac{\partial \mathcal{L}}{\partial
    \mathbf{x}_{\mathrm{flat}}[c \cdot H' \cdot W' + h' \cdot W' + w']}
\end{equation}

\subsubsection{Gradiente a través del MAX Pooling}

El gradiente fluye únicamente hacia la posición ganadora registrada en la máscara;
todas las demás posiciones de la ventana reciben gradiente cero:
\begin{equation}
  \frac{\partial \mathcal{L}}{\partial \mathbf{A}[c, h, w]}
  = \begin{cases}
    \dfrac{\partial \mathcal{L}}{\partial \mathbf{M}[c, h', w']}
      & \text{si } (h, w) \text{ fue el máximo de la ventana } (h',w') \\[8pt]
    0 & \text{en caso contrario}
  \end{cases}
\end{equation}

\subsubsection{Gradiente a través de la Convolución}

Sea $\boldsymbol{\delta} = \frac{\partial \mathcal{L}}{\partial \mathbf{A}}
\odot \mathbf{1}[\mathbf{Z} > 0]$ el delta tras deshacer la ReLU.
Los gradientes respecto a los kernels y al input son:

\begin{align}
  \frac{\partial \mathcal{L}}{\partial \mathbf{K}[f,c,i,j]}
  &= \sum_{h'}\sum_{w'} \boldsymbol{\delta}[f, h', w'] \cdot
     \mathbf{X}_{\mathrm{pad}}[c,\; h' \cdot s + i,\; w' \cdot s + j]
  \label{eq:grad_kernel} \\[6pt]
  \frac{\partial \mathcal{L}}{\partial \mathbf{X}[c, h, w]}
  &= \sum_{f}\sum_{i}\sum_{j} \boldsymbol{\delta}[f, h', w'] \cdot \mathbf{K}[f, c, i, j]
  \label{eq:grad_input_conv}
\end{align}

donde en \eqref{eq:grad_input_conv} la suma recorre todas las ventanas $(h', w')$ cuyo
campo receptivo cubre la posición $(h, w)$. Esta operación es equivalente a una convolución
\textit{full} de $\boldsymbol{\delta}$ con el kernel volteado $180°$.

\subsubsection{Flujo End-to-End del Gradiente}

El pipeline completo de retropropagación sigue el orden inverso al forward:
\begin{equation}
  \underbrace{\frac{\partial \mathcal{L}}{\partial \mathbf{x}_{\mathrm{flat}}}}_{\text{GPU}}
  \;\xrightarrow{\;\mathrm{Flatten}^{-1}\;}
  \frac{\partial \mathcal{L}}{\partial \mathbf{M}}
  \;\xrightarrow{\;\mathrm{Pool}_2^{-1}\;}
  \frac{\partial \mathcal{L}}{\partial \mathbf{A}_2}
  \;\xrightarrow{\;\mathrm{Conv}_2^{-1}\;}
  \frac{\partial \mathcal{L}}{\partial \mathbf{A}_1}
  \;\xrightarrow{\;\mathrm{Pool}_1^{-1}\;}
  \frac{\partial \mathcal{L}}{\partial \mathbf{Z}_1}
  \;\xrightarrow{\;\mathrm{Conv}_1^{-1}\;}
  \frac{\partial \mathcal{L}}{\partial \mathbf{X}}
  \label{eq:gradient_flow}
\end{equation}

% ──────────────────────────────────────────────────────────────
\subsection{Programación Paralela en GPU con CUDA}

\subsubsection{Modelo de Ejecución SIMT}

CUDA organiza la ejecución bajo el modelo \textit{Single Instruction, Multiple Threads}
(SIMT) en una jerarquía de tres niveles:

\begin{itemize}
  \item \textbf{Thread:} unidad mínima de ejecución. Cada thread ejecuta el mismo código
    del kernel pero opera sobre datos distintos.
  \item \textbf{Block:} grupo de threads que comparten \textit{shared memory} y pueden
    sincronizarse con la barrera \texttt{\_\_syncthreads()}.
  \item \textbf{Grid:} colección de bloques que ejecuta el mismo kernel. Los bloques
    son independientes entre sí y el scheduler de la GPU los distribuye libremente
    entre los multiprocesadores (SM).
\end{itemize}

El índice global de un thread en un grid unidimensional se calcula como:
\begin{equation}
  \texttt{idx} = \texttt{blockIdx.x} \times \texttt{blockDim.x} + \texttt{threadIdx.x}
  \label{eq:cuda_idx}
\end{equation}

Los kernels que operan sobre matrices 2D utilizan grids bidimensionales con índices
\texttt{r} (fila) y \texttt{c} (columna) para que cada thread actualice un único
elemento de la matriz.

\subsubsection{Jerarquía de Memoria GPU}

\begin{center}
\begin{tabular}{llll}
\toprule
\textbf{Tipo}      & \textbf{Alcance}   & \textbf{Latencia aprox.} & \textbf{Capacidad aprox.} \\
\midrule
Registros          & Por thread         & 1 ciclo                  & $\sim$64 KB/SM \\
Shared memory      & Por bloque         & $\sim$5 ciclos           & 48--96 KB/bloque \\
Caché L1/L2        & Por SM / chip      & $\sim$30 ciclos          & Automático \\
Memoria global     & Todo el device     & $\sim$800 ciclos         & Varios GB \\
\bottomrule
\end{tabular}
\end{center}

La \textbf{shared memory} es clave para operaciones de reducción paralela como el
cálculo del softmax, donde múltiples threads deben cooperar para encontrar el máximo
y la suma del mismo vector.

\subsubsection{cuBLAS: Multiplicación Matriz-Vector}

La operación central del forward del MLP en cada capa $l$ es:
\begin{equation}
  \mathbf{z}^{[l]} = \bigl(W^{[l]}\bigr)^T \tilde{\mathbf{a}}^{[l-1]}
  \label{eq:matvec}
\end{equation}

Esta operación se delega a \texttt{cublasSgemv}, la rutina cuBLAS para producto
matriz-vector en simple precisión. Como cuBLAS asume almacenamiento en columna
(\textit{column-major}) y los pesos se almacenan en fila (\textit{row-major}), el
flag \texttt{CUBLAS\_OP\_N} sobre la vista transpuesta produce el resultado correcto.

% ============================================================
\section{Arquitectura del Pipeline}
% ============================================================

\subsection{Visión General}

La red completa divide el cómputo entre CPU y GPU:

\begin{center}
\begin{tabular}{llll}
\toprule
\textbf{Componente}   & \textbf{Dispositivo} & \textbf{Precisión}      & \textbf{Parámetros (MNIST)} \\
\midrule
Conv1 (1→32, 3×3)     & CPU                  & \texttt{double} 64-bit  & $1\times32\times3\times3 + 32 = 320$ \\
Pool1 (MAX 2×2)       & CPU                  & \texttt{double} 64-bit  & 0 \\
Conv2 (32→64, 3×3)    & CPU                  & \texttt{double} 64-bit  & $32\times64\times3\times3 + 64 = 18\,496$ \\
Pool2 (MAX 2×2)       & CPU                  & \texttt{double} 64-bit  & 0 \\
Flatten               & CPU                  & \texttt{double} 64-bit  & 0 \\
MLP (3136→256→128→10) & GPU                  & \texttt{float}  32-bit  & $\sim$836\,K \\
\bottomrule
\end{tabular}
\end{center}

\subsection{Flujo Forward para MNIST ($28\times28$, 1 canal)}

\begin{equation*}
[1][28][28]
\xrightarrow{\substack{\text{Conv}(1\to32)\\3\times3,\;\text{SAME}}}
[32][28][28]
\xrightarrow{\text{Pool}(2\times2)}
[32][14][14]
\xrightarrow{\substack{\text{Conv}(32\to64)\\3\times3,\;\text{SAME}}}
[64][14][14]
\xrightarrow{\text{Pool}(2\times2)}
[64][7][7]
\xrightarrow{\text{Flatten}}
\mathbf{x}\in\R^{3136}
\xrightarrow{\text{MLP}}
\hat{y}\in\R^{10}
\end{equation*}

\subsection{Flujo Forward para HASYv2 ($32\times32$, 1 canal)}

\begin{equation*}
[1][32][32]
\xrightarrow{\text{Conv}(1\to32)}
[32][32][32]
\xrightarrow{\text{Pool}}
[32][16][16]
\xrightarrow{\text{Conv}(32\to64)}
[64][16][16]
\xrightarrow{\text{Pool}}
[64][8][8]
\xrightarrow{\text{Flatten}}
\mathbf{x}\in\R^{4096}
\xrightarrow{\text{MLP}}
\hat{y}\in\R^{369}
\end{equation*}

\subsection{Interfaz CPU--GPU: Conversión de Precisión}

En la frontera entre CNN y MLP existe una conversión de precisión explícita:

\begin{itemize}
  \item \textbf{Forward (CPU→GPU):} el vector \texttt{double} de Flatten se convierte a
    \texttt{float} y se transfiere a memoria GPU con \texttt{cudaMemcpy
    (HostToDevice)}.
  \item \textbf{Backward (GPU→CPU):} el vector gradiente \texttt{float} devuelto por el
    MLP se convierte a \texttt{double} antes de pasarse a \texttt{flatten.backward()}.
\end{itemize}

Esta elección preserva la precisión aritmética en los gradientes de los kernels
convolucionales (pequeños, sensibles al redondeo) sin sacrificar la velocidad en las
operaciones matriciales del MLP, donde \texttt{float} ofrece mayor throughput en GPU.

% ============================================================
\section{Implementación}
% ============================================================

\subsection{Estructura de Archivos}

\begin{center}
\begin{tabular}{ll}
\toprule
\textbf{Archivo}                & \textbf{Responsabilidad} \\
\midrule
\texttt{conv\_layer.h/.cpp}     & Capa convolucional: forward y backward en CPU \\
\texttt{pooling\_layer.h/.cpp}  & Capa de pooling: forward y backward en CPU \\
\texttt{flatten\_layer.h/.cpp}  & Aplanado: forward y backward en CPU \\
\texttt{mlp\_cuda.h}            & MLP con kernels CUDA, cuBLAS, Adam/SGD en GPU \\
\texttt{dataset.cpp}            & Carga de MNIST (formato IDX) y HASYv2 (PNG) \\
\texttt{main\_cnn\_cuda.cu}     & Pipeline, entrenamiento, reporte HTML \\
\texttt{linalg.h}               & Tipos Eigen (\texttt{Vector}, \texttt{Matrix}) \\
\texttt{stb\_image.h}           & Carga de PNG/JPG (biblioteca header-only) \\
\bottomrule
\end{tabular}
\end{center}

\subsection{ConvLayer: Forward y Backward}

\subsubsection{Forward}

El método \texttt{forward} sigue cuatro pasos en secuencia:

\begin{enumerate}[noitemsep]
  \item \textbf{Padding:} aplica zero-padding según \eqref{eq:same_pad}.
  \item \textbf{Convolución:} evalúa \eqref{eq:conv} para cada filtro $f$ y cada
    posición $(h', w')$ de la salida.
  \item \textbf{Caché:} guarda \texttt{input\_cache\_}, \texttt{padded\_cache\_}
    y \texttt{pre\_relu\_cache\_} ($\mathbf{Z}$) para el backward.
  \item \textbf{ReLU:} aplica $\max(0, \cdot)$ elemento a elemento sobre $\mathbf{Z}$.
\end{enumerate}

\subsubsection{Backward}

El método \texttt{backward} recibe $\partial\mathcal{L}/\partial\mathbf{A}$ y:

\begin{enumerate}[noitemsep]
  \item Deshace la ReLU multiplicando por su derivada: $\boldsymbol{\delta} =
    \frac{\partial\mathcal{L}}{\partial\mathbf{A}} \odot \mathbf{1}[\mathbf{Z}>0]$.
  \item Acumula el gradiente de los kernels (eq.~\eqref{eq:grad_kernel}) y actualiza
    $\mathbf{K} \mathrel{-}= \eta \cdot \nabla_\mathbf{K}\mathcal{L}$.
  \item Actualiza los sesgos: $b_f \mathrel{-}= \eta \cdot
    \sum_{h',w'} \boldsymbol{\delta}[f, h', w']$.
  \item Calcula y retorna $\partial\mathcal{L}/\partial\mathbf{X}$ mediante
    \eqref{eq:grad_input_conv}.
\end{enumerate}

\subsection{PoolingLayer: Forward con Máscara y Backward}

\subsubsection{Forward: guardado de posiciones ganadoras}

\begin{lstlisting}[caption={MAX Pooling forward con máscara de posiciones ganadoras}]
for (int c = 0; c < C; c++)
  for (int h2 = 0; h2 < Hout; h2++)
    for (int w2 = 0; w2 < Wout; w2++) {
      int best_idx = 0;
      double best_val = -1e30;
      for (int pi = 0; pi < pool_size_; pi++)
        for (int pj = 0; pj < pool_size_; pj++) {
          double v = input[c][h2*s+pi][w2*s+pj];
          if (v > best_val) {
            best_val = v;
            best_idx = pi * pool_size_ + pj;
          }
        }
      output[c][h2][w2] = best_val;
      mask_[c][h2][w2]  = best_idx; // indice del ganador
    }
\end{lstlisting}

\subsubsection{Backward: enrutamiento al ganador}

El gradiente se enruta exclusivamente a la posición que produjo el máximo:

\begin{lstlisting}[caption={MAX Pooling backward — solo el ganador recibe gradiente}]
for (int c = 0; c < C; c++)
  for (int h2 = 0; h2 < Hout; h2++)
    for (int w2 = 0; w2 < Wout; w2++) {
      int idx = mask_[c][h2][w2];
      int pi  = idx / pool_size_;
      int pj  = idx % pool_size_;
      // solo la posicion ganadora acumula el gradiente
      grad_in[c][h2*s+pi][w2*s+pj] += grad_out[c][h2][w2];
    }
\end{lstlisting}

\subsection{MLP\_CUDA}

\subsubsection{Gestión de Memoria GPU: Patrón RAII con \texttt{GpuBuf}}

La estructura \texttt{GpuBuf} encapsula un puntero a memoria GPU bajo el patrón RAII
(\textit{Resource Acquisition Is Initialization}): la memoria se reserva en el
constructor y se libera automáticamente en el destructor, eliminando las fugas de
memoria de GPU:

\begin{lstlisting}[caption={GpuBuf: RAII para memoria GPU}]
struct GpuBuf {
    float* ptr = nullptr;
    int    n   = 0;

    explicit GpuBuf(int size) : n(size) {
        cudaMalloc(&ptr, size * sizeof(float));
        cudaMemset(ptr, 0, size * sizeof(float));
    }
    ~GpuBuf() { if (ptr) cudaFree(ptr); }

    // Copia prohibida: evita doble cudaFree sobre el mismo puntero
    GpuBuf(const GpuBuf&)            = delete;
    GpuBuf& operator=(const GpuBuf&) = delete;

    // Solo movimiento: transfiere la propiedad del puntero
    GpuBuf(GpuBuf&& o) noexcept : ptr(o.ptr), n(o.n)
        { o.ptr = nullptr; o.n = 0; }
};
\end{lstlisting}

\subsubsection{Kernels CUDA Implementados}

\begin{table}[H]
\centering
\caption{Kernels CUDA del módulo \texttt{MLP\_CUDA}}
\label{tab:kernels}
\begin{tabular}{lp{9cm}}
\toprule
\textbf{Kernel}                   & \textbf{Función} \\
\midrule
\texttt{kernel\_bias}             & Prepende el nodo de bias ($a_0 = 1$) al
                                    vector de activaciones para fusionar bias
                                    y pesos en una sola operación matricial. \\
\texttt{kernel\_relu}             & Aplica $\max(0, x)$ elemento a elemento
                                    (forward de capas ocultas). \\
\texttt{kernel\_relu\_d}          & Calcula $\mathbf{1}[x > 0]$, la derivada
                                    de ReLU usada en el backward. \\
\texttt{kernel\_leaky\_relu}      & Aplica $x > 0 \;?\; x : \alpha x$
                                    con $\alpha = 0.01$. \\
\texttt{kernel\_sigmoid}          & Aplica $1/(1+e^{-x})$ (activación alternativa). \\
\texttt{kernel\_softmax}          & Softmax numéricamente estable con reducción
                                    paralela en shared memory. Un único bloque. \\
\texttt{kernel\_mul\_ew}          & Producto elemento a elemento $z_i = x_i y_i$;
                                    usado en backward para $\boldsymbol{\delta}
                                    \odot f'(\mathbf{z})$. \\
\texttt{kernel\_matvec\_row}      & Producto matriz-vector en row-major;
                                    usado en backward de capas ocultas. \\
\texttt{kernel\_ce\_deriv}        & Gradiente combinado Softmax+Cross-Entropy:
                                    $\delta_k = \hat{y}_k - y_k$. \\
\texttt{kernel\_sgd\_update}      & SGD: $W_{rc} \mathrel{-}= \eta\,a_r\,\delta_c$. \\
\texttt{kernel\_adam\_update}     & Adam con momentos $m$, $v$ y corrección de
                                    sesgo (uno por peso). \\
\texttt{kernel\_strip\_bias\_row} & Extrae la submatriz de $W$ sin la fila del bias,
                                    necesaria para propagar el gradiente hacia la CNN. \\
\bottomrule
\end{tabular}
\end{table}

\subsubsection{Softmax Numéricamente Estable con Reducción Paralela}

El kernel \texttt{kernel\_softmax} implementa la forma estable de la ecuación
\eqref{eq:conv} para la capa de salida:
\begin{equation}
  s_k = \frac{e^{z_k - \max_j z_j}}{\displaystyle\sum_{j=0}^{K-1} e^{z_j - \max_j z_j}}
  \label{eq:softmax}
\end{equation}

Restar el máximo antes de calcular las exponenciales previene el desbordamiento
(\textit{overflow}) para activaciones grandes sin cambiar el resultado matemático
(numerador y denominador se dividen por el mismo factor $e^{\max z_j}$).

Para encontrar el máximo y la suma de forma eficiente se emplea una
\textbf{reducción paralela en árbol binario} sobre shared memory. En cada paso los
threads activos se reducen a la mitad:

\begin{lstlisting}[caption={Reducción paralela del máximo en shared memory}]
extern __shared__ float sdata[];
int tid = threadIdx.x;

// Cada thread calcula el maximo parcial de su porcion
float maxVal = -1e30f;
for (int i = tid; i < n; i += blockDim.x)
    maxVal = fmaxf(maxVal, x[i]);
sdata[tid] = maxVal;
__syncthreads();

// Reduccion en arbol: O(log P) pasos con P = blockDim.x
for (int s = blockDim.x / 2; s > 0; s >>= 1) {
    if (tid < s) sdata[tid] = fmaxf(sdata[tid], sdata[tid + s]);
    __syncthreads(); // barrera: todos los threads del bloque esperan
}
maxVal = sdata[0]; // maximo global disponible para todos los threads
\end{lstlisting}

La barrera \texttt{\_\_syncthreads()} garantiza que ningún thread avance al siguiente
nivel del árbol antes de que todos hayan escrito su valor en \texttt{sdata}. La
complejidad se reduce de $O(K)$ secuencial a $O(\log_2 P)$ pasos paralelos.

\subsubsection{Gradiente Combinado Softmax + Cross-Entropy}

Una ventaja algebraica clave es que la derivada combinada de la pérdida
Cross-Entropy con activación Softmax en la capa de salida se simplifica a:
\begin{equation}
  \boldsymbol{\delta}^{[L]} = \frac{\partial \mathcal{L}}{\partial \mathbf{z}^{[L]}}
  = \hat{\mathbf{y}} - \mathbf{y}
  \label{eq:ce_softmax_grad}
\end{equation}

donde $\hat{\mathbf{y}}$ son las probabilidades predichas y $\mathbf{y}$ es el
one-hot del target. El kernel \texttt{kernel\_ce\_deriv} implementa esta operación
elemento a elemento en GPU, evitando calcular explícitamente el jacobiano de softmax.

\subsubsection{Optimizador Adam en GPU}

Adam (Adaptive Moment Estimation) adapta la tasa de aprendizaje por parámetro
manteniendo dos momentos exponencialmente ponderados:
\begin{align}
  m_t &= \beta_1 m_{t-1} + (1-\beta_1)\, g_t
        & \text{(media móvil del gradiente)}
  \label{eq:adam_m}\\
  v_t &= \beta_2 v_{t-1} + (1-\beta_2)\, g_t^2
        & \text{(media móvil del cuadrado del gradiente)}
  \label{eq:adam_v}
\end{align}

Las correcciones de sesgo compensan que $m_0 = v_0 = 0$ infla las estimaciones
iniciales:
\begin{equation}
  \hat{m}_t = \frac{m_t}{1-\beta_1^t}, \qquad
  \hat{v}_t = \frac{v_t}{1-\beta_2^t}
  \label{eq:adam_bc}
\end{equation}

La actualización final divide el gradiente corregido por su desviación estándar
estimada, logrando tasas de aprendizaje adaptativas:
\begin{equation}
  W \;\leftarrow\; W - \eta\,\frac{\hat{m}_t}{\sqrt{\hat{v}_t} + \varepsilon}
  \label{eq:adam_update}
\end{equation}

con $\beta_1 = 0.9$, $\beta_2 = 0.999$, $\varepsilon = 10^{-8}$.
Cada thread CUDA actualiza un peso de forma completamente independiente:

\begin{lstlisting}[caption={Kernel Adam: un thread actualiza un peso independiente}]
__global__ void kernel_adam_update(
    float* W, float* mW, float* vW,
    const float* a, const float* delta,
    float lr, float beta1, float beta2, float eps,
    float bc1, float bc2, int rows, int cols)
{
    int r = blockIdx.x * blockDim.x + threadIdx.x; // fila
    int c = blockIdx.y * blockDim.y + threadIdx.y; // columna
    if (r < rows && c < cols) {
        int   idx = r * cols + c;
        float g   = a[r] * delta[c];                     // gradiente crudo
        mW[idx]   = beta1*mW[idx] + (1.0f-beta1)*g;      // eq. (11)
        vW[idx]   = beta2*vW[idx] + (1.0f-beta2)*g*g;    // eq. (12)
        // bc1 = 1/(1-beta1^t), bc2 = 1/(1-beta2^t)
        W[idx]   -= lr * (mW[idx]*bc1) /
                    (sqrtf(vW[idx]*bc2) + eps);           // eq. (14)
    }
}
\end{lstlisting}

Los factores de corrección de sesgo \texttt{bc1} y \texttt{bc2} se calculan una sola
vez en CPU antes de lanzar el kernel, eliminando la necesidad de calcular potencias
dentro de la GPU.

\subsubsection{Gradiente del Input del MLP hacia la CNN}

El gradiente de la pérdida respecto al input del MLP (salida del Flatten) es:
\begin{equation}
  \frac{\partial \mathcal{L}}{\partial \mathbf{x}_{\mathrm{flat}}}
  = W^{[0]}_{\mathrm{sin\,bias}} \cdot \boldsymbol{\delta}^{[0]}
  \label{eq:grad_input_mlp}
\end{equation}

donde $W^{[0]}_{\mathrm{sin\,bias}}$ es la submatriz de $W^{[0]}$ excluyendo la fila
del nodo bias (extraída por \texttt{kernel\_strip\_bias\_row}). El kernel
\texttt{kernel\_matvec\_row} realiza este producto en GPU y el resultado se transfiere
a CPU en \texttt{float} para continuar el backward por la CNN.

\subsubsection{Interfaz Paso a Paso para el Pipeline CNN}

Para integrarse con la CNN, el MLP expone dos métodos públicos en lugar de un ciclo
de entrenamiento autónomo:

\begin{itemize}
  \item \texttt{stepForward(x)}: recibe \texttt{vector<float>} desde el Flatten, lo
    sube a GPU, propaga por todas las capas y devuelve las probabilidades softmax en
    CPU.
  \item \texttt{stepBackward(target, lr)}: recibe el one-hot del target,
    ejecuta el backward completo actualizando los pesos del MLP, y devuelve el
    gradiente respecto al input (vector \texttt{float}) que se pasa a la CNN.
\end{itemize}

Los buffers GPU (\texttt{step\_d\_acts\_}, \texttt{step\_d\_nets\_}, etc.) se
inicializan una sola vez de forma \textit{lazy} en la primera llamada y se reutilizan
en cada iteración, eliminando el costo de \texttt{cudaMalloc} por muestra.

% ============================================================
\section{Datasets}
% ============================================================

\subsection{MNIST}

MNIST es el conjunto de referencia estándar para clasificación de dígitos escritos a
mano. Contiene 60\,000 imágenes de entrenamiento y 10\,000 de prueba, cada una de
$28 \times 28$ píxeles en escala de grises, distribuidas en 10 clases (dígitos 0--9).

Las imágenes se almacenan en el \textbf{formato IDX binario}: cabecera de cuatro
enteros de 32 bits en big-endian (\texttt{magic}, \texttt{n\_samples}, \texttt{rows},
\texttt{cols}) seguida de los bytes de píxel. Cada píxel se normaliza al rango $[0,1]$
dividiendo por 255 antes de alimentar la red.

En la implementación se utilizan 3\,000 muestras de entrenamiento y 500 de prueba para
agilizar el ciclo de desarrollo sin perder representatividad estadística.

\subsection{HASYv2}

HASYv2 es un dataset de símbolos matemáticos escritos a mano que contiene
aproximadamente 168\,000 imágenes de $32 \times 32$ píxeles en escala de grises,
distribuidas en \textbf{369 clases} (caracteres \LaTeX\ como $\alpha$, $\Sigma$,
$\int$, $\infty$, $\nabla$, etc.). El dataset está organizado en 10 folds para
validación cruzada; la implementación utiliza el fold~1.

Las imágenes se almacenan como archivos PNG referenciados desde archivos CSV con el
mapeo \texttt{path, symbol\_id, latex, user\_id}. La clase \texttt{HASYLoader}
reconstruye el mapeo \texttt{symbol\_id}~$\to$~índice de clase a partir del archivo
\texttt{symbols.csv}.

\subsection{Jerarquía Polimórfica: \texttt{DatasetLoader}}

La carga de datos se organiza con una interfaz abstracta:

\begin{lstlisting}[caption={Interfaz abstracta DatasetLoader}]
class DatasetLoader {
public:
    virtual ~DatasetLoader() = default;
    virtual DatasetInfo cargar() = 0;
};
\end{lstlisting}

Esto permite seleccionar el dataset por argumento de línea de comandos sin
modificar el pipeline de entrenamiento. La arquitectura del MLP se ajusta
automáticamente: para datasets con más de 100 clases (HASYv2) se usa una capa
oculta adicional de mayor dimensión.

\subsection{Mezcla Aleatoria por Época}

El conjunto de entrenamiento se mezcla con \texttt{std::shuffle} antes de cada época
usando un generador Mersenne Twister con semilla fija (\texttt{mt19937(42)}):

\begin{lstlisting}[caption={Shuffle del orden de muestras por época}]
vector<int> order(N);
iota(order.begin(), order.end(), 0);
shuffle(order.begin(), order.end(), local_rng);
\end{lstlisting}

La mezcla rompe las correlaciones entre muestras consecutivas, reduce el sesgo del
gradiente estocástico y mejora la generalización del modelo. La semilla fija garantiza
reproducibilidad entre ejecuciones.

% ============================================================
\section{Resultados}
% ============================================================

\subsection{Configuración Experimental}

\begin{center}
\begin{tabular}{lll}
\toprule
\textbf{Parámetro}        & \textbf{MNIST}                         & \textbf{HASYv2} \\
\midrule
Imágenes entrenamiento     & 3\,000                                 & 5\,000 \\
Imágenes prueba            & 500                                    & 1\,000 \\
Tamaño imagen              & $28 \times 28$                         & $32 \times 32$ \\
Clases                     & 10                                     & 369 \\
Arquitectura CNN           & Conv(1$\to$32) + Conv(32$\to$64)       & Conv(1$\to$32) + Conv(32$\to$64) \\
Arquitectura MLP           & 3136$\to$256$\to$128$\to$10            & 4096$\to$512$\to$256$\to$369 \\
Activación oculta          & ReLU                                   & ReLU \\
Activación salida          & Softmax                                & Softmax \\
Inicialización             & He                                     & He \\
Función de pérdida         & Cross-Entropy                          & Cross-Entropy \\
Optimizador MLP            & Adam ($\eta_{\text{MLP}}=0.001$)       & Adam ($\eta_{\text{MLP}}=0.001$) \\
Optimizador CNN            & SGD  ($\eta_{\text{CNN}}=0.001$)       & SGD  ($\eta_{\text{CNN}}=0.001$) \\
Épocas                     & 50                                     & 50 \\
Padding                    & SAME                                   & SAME \\
Pooling                    & MAX $2\times2$                         & MAX $2\times2$ \\
\bottomrule
\end{tabular}
\end{center}

\subsection{Curvas de Aprendizaje}

\todo{Insertar gráficas de pérdida (train loss) y precisión (train accuracy) por época
para MNIST y HASYv2. Generadas por el reporte HTML del programa.}

\subsection{Matriz de Confusión y Precisión por Clase}

\todo{Insertar matriz de confusión y tabla de precisión por clase para MNIST (test).
Para HASYv2 (369 clases) incluir solo la tabla de precisión por clase.}

\subsection{Comparativa: MLP Puro vs CNN+MLP}

\begin{table}[H]
\centering
\caption{Comparativa MLP puro vs CNN+MLP en MNIST (500 muestras test)}
\label{tab:comparativa}
\begin{tabular}{lll}
\toprule
\textbf{Característica}        & \textbf{MLP puro}              & \textbf{CNN+MLP} \\
\midrule
Input a la red                 & Píxeles crudos ($784$ dim)     & Mapa de características ($3136$ dim) \\
Arquitectura                   & $784\to256\to128\to10$         & CNN + $3136\to256\to128\to10$ \\
Parámetros totales             & $\sim$235\,K                   & $\sim$855\,K \\
Extracción de características  & Ninguna                        & Automática (filtros aprendidos) \\
Compartición de parámetros     & No                             & Sí (mismo filtro en todas las posiciones) \\
Invarianza traslacional        & No                             & Parcial (via MAX pooling) \\
Optimizador                    & Adam                           & Adam (MLP) + SGD (CNN) \\
Precisión test MNIST           & $\sim$89.6\%                   & \todo{resultado} \\
Tiempo por época (MNIST)       & \todo{resultado}               & \todo{resultado} \\
\bottomrule
\end{tabular}
\end{table}

La ventaja principal de la CNN no reside únicamente en la precisión final, sino en su
capacidad de aprender \textbf{representaciones jerárquicas}: la primera capa
convolucional detecta bordes y texturas locales, la segunda combina esas características
en patrones más complejos. El MLP opera sobre estas representaciones compactas en lugar
de sobre los píxeles brutos, requiriendo menos capacidad para alcanzar igual o mayor
precisión.

% ============================================================
\section{Conclusiones}
% ============================================================

\begin{enumerate}
  \item \textbf{Compartición de parámetros:} Un filtro de $3\times3$ aplicado sobre
    un mapa $28\times28$ utiliza solo 9 pesos para procesar 784 posiciones, frente a
    los $784\times N_1$ pesos de una capa densa equivalente. Esto reduce el sobreajuste
    y mejora la generalización a datos no vistos.

  \item \textbf{Retropropagación end-to-end:} La implementación demuestra que el
    gradiente puede fluir de forma continua desde la capa de salida del MLP hasta los
    kernels de la primera capa convolucional, siguiendo el camino descrito en
    \eqref{eq:gradient_flow}, sin necesidad de un framework de diferenciación automática.

  \item \textbf{División CPU/GPU:} Mantener las capas CNN en CPU (\texttt{double}) y
    el MLP en GPU (\texttt{float}) es un compromiso práctico que facilita la
    implementación correcta del backward convolucional mientras aprovecha el paralelismo
    masivo de CUDA para las operaciones densas del MLP.

  \item \textbf{Softmax estable con reducción paralela:} Restar el máximo antes de
    las exponenciales garantiza estabilidad numérica. La reducción en shared memory con
    \texttt{\_\_syncthreads()} reduce la complejidad de $O(K)$ a $O(\log_2 P)$ pasos.

  \item \textbf{Gradiente simplificado Softmax+CE:} La combinación algebraica de
    Softmax y Cross-Entropy reduce el delta de la capa de salida a
    $\boldsymbol{\delta}^{[L]} = \hat{\mathbf{y}} - \mathbf{y}$ (ecuación
    \eqref{eq:ce_softmax_grad}), eliminando el cálculo explícito del jacobiano de
    softmax.

  \item \textbf{Adam vs SGD:} Adam converge más rápido para los parámetros densos del
    MLP (muchos pesos, gradientes ruidosos). SGD es suficiente para los kernels
    convolucionales (pocos pesos, gradientes más estables), justificando el uso de
    optimizadores distintos en cada parte del pipeline.

  \item \textbf{Patrón RAII con \texttt{GpuBuf}:} Encapsular la memoria GPU con
    constructor/destructor y semántica de movimiento garantiza que no haya fugas de
    memoria ni doble \texttt{cudaFree}, aplicando en CUDA los principios de gestión
    segura de recursos de C++ moderno.
\end{enumerate}

% ============================================================
\section*{Compilación y Ejecución}
% ============================================================

\subsubsection*{Dependencias}

\begin{lstlisting}[numbers=none]
# CUDA Toolkit >= 11 (incluye nvcc y cuBLAS)
sudo apt install cuda-toolkit-12-x

# Eigen (algebra lineal en CPU)
sudo apt install libeigen3-dev

# stb_image.h ya esta incluido en el repositorio
\end{lstlisting}

\subsubsection*{Compilación}

\begin{lstlisting}[numbers=none]
nvcc -O2 -std=c++17 -lcublas -o cnn_cuda main_cnn_cuda.cu
\end{lstlisting}

\subsubsection*{Ejecución}

\begin{lstlisting}[numbers=none]
./cnn_cuda mnist    # clasificacion de digitos 0-9
./cnn_cuda hasy     # clasificacion de simbolos matematicos (369 clases)
\end{lstlisting}

Al finalizar el entrenamiento, el programa imprime en consola el resumen de resultados,
la matriz de confusión (para datasets con hasta 30 clases) y el historial completo de
pérdida y precisión por época. También genera automáticamente un reporte HTML
interactivo con curvas de aprendizaje (Chart.js), predicciones visuales con canvas y
matriz de confusión codificada por color.

\end{document}
